\section{Etapa 8 --- Auditor\'ia pre-campo}
\label{sec:et8}
\periodo{11 --- 12 de julio de 2026 \quad$\cdot$\quad 2 d\'ias, 33 \emph{commits}}

Dos d\'ias enteros dedicados exclusivamente a \emph{tratar de romper el
sistema antes de que lo rompiera el campo}. Es, en retrospectiva, la etapa
mejor invertida del proyecto: encontr\'o tres defectos que habr\'ian arruinado
campa\~nas completas y que no se manifestaban en las pruebas cortas de
laboratorio.

\begin{figure}[ht]
\centering
\includegraphics[width=\textwidth]{flujo_adquisicion_masw.png}
\caption{Flujo que la auditor\'ia busc\'o preservar de extremo a extremo. Se distinguen
integridad de adquisici\'on, reconstrucci\'on y validaci\'on geot\'ecnica porque no son el
mismo criterio de aceptaci\'on.}
\end{figure}

\subsection{Por qu\'e hac\'ia falta}

Todo lo que se hab\'ia validado hasta ese momento eran capturas cortas. Salir a
campo implica sesiones largas, reconexiones, cortes de energ\'ia, operaci\'on por
interfaz web bajo sol y con una sola oportunidad por sitio. Los modos de falla
que importan en esas condiciones son distintos de los que aparecen en el
banco.

\subsection{La jerarqu\'ia de memoria, validada}

Primero hubo que confirmar que el camino de datos completo funcionaba en las
tres escalas previstas:

\begin{center}\small
\begin{tabularx}{\textwidth}{@{}l X l@{}}
\toprule
\textbf{Nivel} & \textbf{D\'onde vive la captura} & \textbf{Alcance} \\
\midrule
1 & RAM del ESP32 & capturas cortas \\
2 & Encadenado hacia el PSoC & capturas medias \\
3 & \textbf{Tarjeta SD, gestionada por el PSoC} (FatFs, \texttt{GEOLAST.BIN}, comandos \texttt{0xBC}--\texttt{0xBF}) & capturas largas \\
\bottomrule
\end{tabularx}
\end{center}

\seleccionado{La SD vive en el \textbf{PSoC}, no en el ESP. Es coherente con la
divisi\'on de responsabilidades del sistema: el PSoC posee la se\~nal y su
temporizado, el ESP transporta. Se agreg\'o el comando de banco
\texttt{sdread N} para leer bancos directamente y poder auditar lo escrito.}

\subsection{La matriz de aceptaci\'on E1--E18}

\begin{center}\small
\begin{tabularx}{\textwidth}{@{}l X l@{}}
\toprule
\textbf{Id} & \textbf{Prueba} & \textbf{Resultado} \\
\midrule
E1 & RAM chica por radio (\texttt{-n 4}) & PASS 120/120 \\
E2 & SD chica por radio (\texttt{-n 600}, \SI{6.9}{\second}) & PASS 18\,000/18\,000 \\
E3 & SD mediana por radio (\texttt{-n 3000}, \SI{34.6}{\second}) & PASS 90\,000/90\,000 \\
E4 & Log paralelo durante E3 & \SI{3.6}{\percent} de lecturas SD sobre \SI{250}{\milli\second} (defecto H3) \\
E5 & \textbf{Aceptaci\'on de 10 minutos} (52\,080 lotes) & \textbf{FAIL --- 0 muestras} \\
E5b/E5c & Repetici\'on con el fix F1 & \textbf{PASS extremo a extremo} \\
E6/E6b & Decimaci\'on 2 y FIR por radio & PASS 3000/3000 \\
E7--E9 & Comandos \texttt{0xBE}, \texttt{0xAD}, USB estricto & PASS (33/33 en USB) \\
E10 & Interfaz en navegador & PASS \\
E12 & \textbf{Reconexi\'on WebSocket cortada durante el volcado} & PASS exacto \\
E13 & STOP en vivo, en RUNNING y en DUMPING & PASS \\
E14 & Operaci\'on completa desde la interfaz & PASS \\
E15 & Captura SD m\'axima, 60\,000 lotes & PASS \\
E16 & Transiciones de rango y decimaci\'on & PASS 49/49 \\
E17 & Ciclo de energ\'ia y recuperaci\'on & PASS 30/30 \\
E18 & \emph{Mini-soak} & PASS 5/5 \\
\bottomrule
\end{tabularx}
\end{center}

\dato{\textbf{Resultado de la prueba de aceptaci\'on final (E5c):}
52\,080/52\,080 lotes escritos en SD y 1\,562\,400/1\,562\,400 muestras
transferidas por radio, sobre una captura continua de 10 minutos. Cero
p\'erdidas. Veredicto registrado: \emph{sistema del banco listo para campo}.}

\subsection{Los tres defectos cr\'iticos}

\subsubsection{F1 --- el \emph{watchdog} destru\'ia las capturas largas}

El defecto que produjo el ``10 minutos $=$ 0 muestras'', y el m\'as instructivo
de todos: \textbf{el mecanismo de protecci\'on era la causa de la falla}.

El \emph{watchdog} \texttt{SYNC\_STALL} exist\'ia para detectar que una captura
se hab\'ia trabado. Su margen era fijo, no proporcional a la duraci\'on de la
captura, as\'i que en sesiones largas ---donde escribir en SD demora
leg\'itimamente--- disparaba solo, y su acci\'on de recuperaci\'on era
\emph{re-disparar} \texttt{START\_NOW}. Cada re-disparo destru\'ia la sesi\'on en
curso y arrancaba otra, que a su vez volv\'ia a disparar el \emph{watchdog}. El
resultado eran diez minutos de trabajo y cero muestras.

\seleccionado{\textbf{Fix:} margen proporcional a la duraci\'on esperada, y
aborto limpio en vez de re-disparo. La lecci\'on transferible es que un
\emph{watchdog} cuya acci\'on correctiva reinicia el trabajo puede convertir una
lentitud tolerable en una p\'erdida total, y que \textbf{esa clase de defecto
solo aparece en pruebas de duraci\'on realista}.}

\subsubsection{F9 --- START romp\'ia el nodo}

\dato{Marcado como cr\'itico: \emph{START romp\'ia el nodo y lo dejaba inservible
hasta reset f\'isico, siempre, con cualquier $N$}. Causa ra\'iz: una condici\'on
de carrera entre el reintento de PRESTART del maestro y un estado ARMED del
PSoC que no se anunciaba. El maestro reintentaba mientras el PSoC ya hab\'ia
transicionado en silencio, y el nodo quedaba en un estado inconsistente del
que no sal\'ia solo.}

Es la contracara del logro de la Etapa~4: el mismo mecanismo de
sincronizaci\'on que resolvi\'o el arranque coordinado ten\'ia un camino de
reintento sin cubrir.

\subsubsection{F5 --- volcado colgado por socket zombi}

La herramienta de captura por WebSocket se colgaba indefinidamente si el
socket quedaba medio abierto: sin datos, sin error, sin fin. Se corrigi\'o con
un \emph{watchdog} de silencio y endurecimiento previo a la fase de datos, y
se valid\'o cortando la conexi\'on a prop\'osito durante el volcado (E12), que
reconect\'o y termin\'o bien.

\subsection{Defectos menores documentados y no corregidos}

Registrarlos vale tanto como corregirlos, porque evita volver a investigarlos:

\begin{center}\small
\begin{tabularx}{\textwidth}{@{}l X l@{}}
\toprule
\textbf{Id} & \textbf{Descripci\'on} & \textbf{Estado} \\
\midrule
H3 & \SI{3.6}{\percent} de lecturas SD superan \SI{250}{\milli\second} (picos de \SIrange{1.2}{1.5}{\second}) & Mitigado: \emph{timeout} \SI{250}{}$\to$\SI{450}{\milli\second}; el reintento siempre salva \\
F2 & Los ACK muestran \texttt{value \& 0xFF} (600 se ve como 88) & Cosm\'etico, documentado \\
W3 & STOP durante el volcado: no hay re-volcado desde la interfaz (los datos quedan en la SD) & Documentado \\
W5 & La vista previa de $N$ usa el campo de entrada y no el ACK & Documentado \\
H4 & \emph{Prefetch} de bloque en el PSoC dar\'ia un volcado \textasciitilde{}$2\times$ m\'as r\'apido & Mejora identificada, post-campo \\
W6/W7 & M\'aximo de RAM en 513 lotes por redondeo hacia arriba; paso del campo de duraci\'on & Corregidos \\
\bottomrule
\end{tabularx}
\end{center}

\subsection{Balance de la etapa}

\begin{center}\small
\begin{tabularx}{\textwidth}{@{}l X@{}}
\toprule
\textbf{Se qued\'o} & Jerarqu\'ia de memoria RAM$\to$encadenado$\to$SD validada en hardware; SD gestionada por el PSoC con FatFs; matriz E1--E18 como procedimiento de aceptaci\'on repetible; suite de \emph{scripts} de regresi\'on (USB, SD m\'axima, transiciones de rango, recuperaci\'on por reset); \emph{watchdog} de silencio en las herramientas \\
\addlinespace[3pt]
\textbf{Se descart\'o} & \emph{Watchdog} con margen fijo; re-disparo de \texttt{START\_NOW} como acci\'on de recuperaci\'on; criterios de prueba que se bloqueaban por contadores benignos (F8) \\
\addlinespace[3pt]
\textbf{Aprendizaje} & Los defectos que importan en campo son de \emph{duraci\'on} y de \emph{reconexi\'on}, y no aparecen en pruebas cortas. Tres de ellos habr\'ian costado campa\~nas enteras \\
\bottomrule
\end{tabularx}
\end{center}
